\documentclass[10pt,a4paper]{amsart}
\usepackage[margin=19mm]{geometry}
\usepackage{amsmath,amssymb,mathtools}
\usepackage[scr=rsfs,cal=euler]{mathalfa}
\usepackage{xfrac}
\usepackage[unicode,hidelinks,bookmarks=false]{hyperref}
\newcommand{\ie}{i.e.\ }
\newcommand{\cf}{cf.\ }
\newcommand{\resp}{respectively\ }
\newcommand{\Subsection}{\subsection}
\providecommand{\nobreakdash}{\nobreak\hbox{-}}
\setlength{\emergencystretch}{2em}
\multlinegap0pt
\usepackage{mathtools}
\usepackage{amssymb}
\usepackage{xcolor}
\usepackage{float}
\usepackage{tikz}
\newtheorem{theorem}{Theorem}
\newtheorem{lemma}{Lemma}
\newcommand{\C}{\mathbb C}%
\newcommand{\Hb}{\mathbb H}%
\newcommand{\N}{\mathbb N}%
\newcommand{\R}{\mathbb R}%
\newcommand{\Sb}{\mathbb S}%
\newcommand{\X}{\mathbb X}%
\title{Regularity and pointwise convergence for dispersive equations on Riemannian symmetric spaces of compact type}
\author{Utsav Dewan, Sanjoy Pusti}
\date{Source: arXiv:2512.09689v3 (CC BY 4.0)}
\begin{document}
\maketitle

\noindent\emph{Source and licence.} Regularity and pointwise convergence for dispersive equations on Riemannian symmetric spaces of compact type. Version arXiv:2512.09689v3 (CC BY 4.0), distributed by arXiv under the Creative Commons Attribution 4.0 International licence, http://creativecommons.org/licenses/by/4.0/, as recorded in the arXiv licence metadata for this exact version and shown on the abstract page https://arxiv.org/abs/2512.09689v3. Full licence notice: \texttt{licenses/arxiv-2512.09689-CC-BY-4.0.txt}.\par\smallskip

\noindent\textit{Editorial scope.} Counted unit: the article's theorem transference\_principle (source0.tex lines 284--286). The article's complete original proof of it is delivered with it (source0.tex lines 722--828, one \texttt{proof} environment of 107 lines, which the article places away from the statement). No mathematics is written, repaired or re-derived: the statement and the proof are the article's own lines, in the article's own notation, and no retained line is substituted or re-flowed.  Retained uncounted as the source-local context the proof needs: lines 486--488; lines 497--514; lines 701--707.  Nothing was omitted from any retained span, so the package carries no omission record.  No retained line contains a citation, so the wrapper carries no bibliography and no bibliography entry.  No retained line contains a figure environment, an image-inclusion command or drawing code, so no image is delivered and the package declares no asset dependency.  Editorial formatting only, in addition: an AMS \texttt{amsart} preamble that restates the article's own declarations and macros used by the retained lines, namely the environments \texttt{theorem}, \texttt{lemma} on separate counters, together with the standard packages those lines need.  The retained environments print in this document as Theorem 1, Lemma 1 in order of appearance, whereas the article prints the same items with the numbering of its own full document.  The complete native source of the article as distributed in the arXiv e-print for this exact version is bundled unchanged as \texttt{sources/190753/source0.tex}.
\begin{equation} \label{spherical_fourier_series}
f=\sum_{n=0}^\infty a_n\:\tilde{Z}_n\:,
\end{equation}

\begin{lemma} \label{Linfty_bound}
\begin{itemize}
\item[(i)] Near the origin, we have
 \begin{equation*}
 \|\tilde{Z}_n\|_{L^\infty\left(\left[0,\frac{\pi}{2}\right]\right)} \lesssim {\left(1+n\right)}^{\frac{d-1}{2}}\:,\:\: n \in \N \cup \{0\}\:.
\end{equation*}  
\item[(ii)] Near the antipodal manifold, we have
\begin{equation*}
\|\tilde{Z}_n\|_{L^\infty\left(\left[\frac{\pi}{2},\pi\right]\right)} \lesssim \begin{cases}
(1+n)^{\frac{d-1}{2}}\:\:\:&\text{if } \X=\Sb^d\:,\\
1\:\:\:&\text{if } \X=P^d(\R)\:,\\
(1+n)^{\frac{1}{2}}\:\:\:&\text{if } \X=P^d(\C)\:,\\
(1+n)^{\frac{3}{2}}\:\:\:&\text{if } \X=P^d(\Hb)\:,\\
(1+n)^{\frac{7}{2}}\:\:\:&\text{if } \X=P^{16}(Cay)\:.
\end{cases}
\end{equation*}
\end{itemize}
\end{lemma}

\begin{lemma} \label{sobolev_comparison}
Let $f \in C^\infty(\X)$ with spherical Fourier series $f=\displaystyle\sum_{n=N}^\infty a_n \tilde{Z}_n$ (as in (\ref{spherical_fourier_series})), for some $N \in \N$. Then for $0 \le \beta_1 < \beta_2 < \infty$, we have
\begin{equation*}
\|f\|_{H^{\beta_1}(X)} \lesssim N^{-(\beta_2-\beta_1)}\:\|f\|_{H^{\beta_2}(X)}\:,
\end{equation*}
with the implicit constant depending only on the intrinsic geometry of $\X$.
\end{lemma}

\begin{theorem} \label{transference_principle}
Let $\psi_1$ and $\psi_2$ be continuous real-valued functions on $[0,\infty)$. If the dispersive equations corresponding to $\psi_1$ and $\psi_2$ are of comparable oscillation, then they are also transferrable.  
\end{theorem} 

\begin{proof}[Proof of Theorem \ref{transference_principle}]
Without loss of generality, let us assume that for some $\alpha_0>0$ and some $p \in [1,\infty]$, the maximal estimate 
\begin{equation} \label{transference_pf_eq1}
\left\|\sup_{0 \le t < 2\pi} \left|e^{-it\psi_1\left(\sqrt{-\Delta}\right)} f \right|\right\|_{L^p(\X)} \lesssim {\|f\|}_{H^\alpha(\X)}\:,
\end{equation}
holds for all $\alpha >\alpha_0$ and all $K$-biinvariant $f \in C^\infty(\X)$. We have to show that the maximal estimate 
\begin{equation} \label{transference_pf_eq2}
\left\|\sup_{0 \le t < 2\pi} \left|e^{-it\psi_2\left(\sqrt{-\Delta}\right)} f \right|\right\|_{L^p(\X)} \lesssim {\|f\|}_{H^\alpha(\X)}\:,
\end{equation}
also holds for all $\alpha >\alpha_0$ and all $K$-biinvariant $f \in C^\infty(\X)$.

By the hypothesis there exist positive constants $R$ and $C$ such that 
\begin{equation} \label{transference_pf_eq3}
\left|\psi_1(r)-\psi_2(r)\right| \le C\:,\:\: r \in (R, \infty)\:.
\end{equation} 
Let us choose and fix $f \in C^\infty(\X)$. Correspondingly, we have the spherical Fourier series (as in (\ref{spherical_fourier_series})),
\begin{equation*}
f = \sum_{n=0}^\infty a_n \tilde{Z}_n\:.
\end{equation*}
Let $m_0$ be the smallest positive integer such that $2^{m_0}> R$. We decompose $f$ as follows,
\begin{equation*} 
f=f_{m_0} + \sum_{m=m_0+1}^\infty f_m\:,
\end{equation*}
where 
\begin{equation*}
f_{m_0}:= \sum_{n=0}^{2^{m_0+1}-1} a_n \tilde{Z}_n \:, \text{ and for } m \ge m_0+1,\:\:\: f_m :=  \sum_{n=2^m}^{2^{m+1}-1} a_n \tilde{Z}_n\:.
\end{equation*}
Invoking the above decomposition of $f$, we get that
\begin{eqnarray} \label{transference_pf_eq4}
&&\left\|\sup_{0 \le t < 2\pi} \left|e^{-it\psi_2\left(\sqrt{-\Delta}\right)} f \right|\right\|_{L^p(\X)}\nonumber\\
 & \le & \left\|\sup_{0 \le t < 2\pi} \left|e^{-it\psi_2\left(\sqrt{-\Delta}\right)} f_{m_0} \right|\right\|_{L^p(\X)} + \sum_{m=m_0+1}^\infty \left\|\sup_{0 \le t < 2\pi} \left|e^{-it\psi_2\left(\sqrt{-\Delta}\right)} f_m \right|\right\|_{L^p(\X)}\:.
\end{eqnarray}
Since $f$ is $K$-biinvariant, so is each of its pieces $f_m$ above, their propagations $e^{-it\psi_j\left(\sqrt{-\Delta}\right)} f_m$, for $j=1,2$ and also the corresponding maximal functions.  

By the $L^\infty$ bound given in (\ref{Linfty_bound}) and the Cauchy-Schwarz inequality we have for all $t \in [0,2\pi)$, all $\theta \in [0,\pi]$ and all $\alpha>\alpha_0$,
\begin{eqnarray*}
\left|e^{-it\psi_2\left(\sqrt{-\Delta}\right)} f_{m_0}(\theta) \right| &=& \left|\sum_{n=0}^{2^{m_0+1}-1} e^{-it\psi_2\left(\lambda_n\right)} a_n \tilde{Z}_n \right| \\
&\le & \sum_{n=0}^{2^{m_0+1}-1} |a_n|\:\left|\tilde{Z}_n(\theta)\right|\\
&\lesssim & \|f_{m_0}\|_{H^{\alpha}(\X)} \\
& \le & \|f\|_{H^{\alpha}(\X)}\:,
\end{eqnarray*}
from which it follows that
\begin{equation} \label{transference_pf_eq5}
\left\|\sup_{0 \le t < 2\pi} \left|e^{-it\psi_2\left(\sqrt{-\Delta}\right)} f_{m_0} \right|\right\|_{L^p(\X)} \lesssim  \|f\|_{H^{\alpha}(\X)}\:.
\end{equation}
Next for each $m \ge m_0+1$, we have for all $\theta \in [0,\pi]$ and all $t \in [0,2\pi)$,
\begin{eqnarray*}
&&\left|e^{-it\psi_1\left(\sqrt{-\Delta}\right)} f_m(\theta)   - e^{-it\psi_2\left(\sqrt{-\Delta}\right)} f_m(\theta) \right| \\
&=& \left|\sum_{n=2^m}^{2^{m+1}-1} \left(e^{-it\psi_1\left(\lambda_n\right)} - e^{-it\psi_2\left(\lambda_n\right)}\right) a_n \tilde{Z}_n(\theta) \right|\\
&=& \left|\sum_{n=2^m}^{2^{m+1}-1} e^{-it\psi_1\left(\lambda_n\right)}\left(1 - e^{-it\left(\psi_2\left(\lambda_n\right)-\psi_1\left(\lambda_n\right)\right)}\right) a_n \tilde{Z}_n(\theta) \right|\:.
\end{eqnarray*}
Now using the smoothness of the exponential in time, we expand the exponential in the Taylor series to get,
\begin{eqnarray*}
&&\left|\sum_{n=2^m}^{2^{m+1}-1} e^{-it\psi_1\left(\lambda_n\right)}\left(1 - e^{-it\left(\psi_2\left(\lambda_n\right)-\psi_1\left(\lambda_n\right)\right)}\right) a_n \tilde{Z}_n(\theta) \right| \\
&=& \left|\sum_{n=2^m}^{2^{m+1}-1} e^{-it\psi_1\left(\lambda_n\right)}\left[1 - \sum_{j=0}^\infty \frac{\left\{-it\left(\psi_2\left(\lambda_n\right)-\psi_1\left(\lambda_n\right)\right)\right\}^j}{j!}\right]
a_n \tilde{Z}_n(\theta) \right|\\
& \le & \sum_{j=1}^\infty \frac{t^j}{j!} \left|\sum_{n=2^m}^{2^{m+1}-1}e^{-it\psi_1\left(\lambda_n\right)} \left[\psi_2(\lambda_n)-\psi_1(\lambda_n)\right]^j a_n\:\tilde{Z}_n(\theta)\right|\:.
\end{eqnarray*}
For each $j \in \N$ and $n \in \left\{2^m, \dots, 2^{m+1}-1\right\}$, setting $a_{j,n}:= \left[\psi_2(\lambda_n)-\psi_1(\lambda_n)\right]^j a_n$, we define,
\begin{equation*}
g_{j,m}:= \sum_{n=2^m}^{2^{m+1}-1} a_{j,n} \tilde{Z}_n\;.
\end{equation*}
Then the computation above yields, for each $m \ge m_0+1$, for all $\theta \in [0,\pi]$ and all $t \in [0,2\pi)$,
\begin{equation*}
\left|e^{-it\psi_1\left(\sqrt{-\Delta}\right)} f_m(\theta)   - e^{-it\psi_2\left(\sqrt{-\Delta}\right)} f_m(\theta) \right| \le  \sum_{j=1}^\infty \frac{t^j}{j!} \left|e^{-it\psi_1\left(\sqrt{-\Delta}\right)} g_{j,m}(\theta)\right|\:,
\end{equation*}
which in turn implies
\begin{eqnarray} \label{transference_pf_eq6}
&&\left\|\sup_{0\le t<2\pi}\left|e^{-it\psi_1\left(\sqrt{-\Delta}\right)} f_m   - e^{-it\psi_2\left(\sqrt{-\Delta}\right)} f_m \right|\right\|_{L^p(\X)} \nonumber\\
&\le & \sum_{j=1}^\infty \frac{1}{j!} \left\| \sup_{0\le t<2\pi} \left|e^{-it\psi_1\left(\sqrt{-\Delta}\right)} g_{j,m}\right|\right\|_{L^p(\X)}\:.
\end{eqnarray}
For any $\varepsilon>0$, set $\alpha_1:=\alpha_0+\frac{\varepsilon}{2}$ and $\alpha_2:=\alpha_0+\varepsilon$. Then applying the assumption (\ref{transference_pf_eq1}) and Lemma \ref{sobolev_comparison}, we get that
\begin{equation*}
\left\| \sup_{0\le t<2\pi} \left|e^{-it\psi_1\left(\sqrt{-\Delta}\right)} g_{j,m}\right|\right\|_{L^p(\X)} \lesssim \|g_{j,m}\|_{H^{\alpha_1}(\X)} \lesssim 2^{-\frac{m\varepsilon}{2}}  \|g_{j,m}\|_{H^{\alpha_2}(\X)}\:,
\end{equation*}
which upon plugging in (\ref{transference_pf_eq6}) yields,
\begin{equation} \label{transference_pf_eq7}
\left\|\sup_{0\le t<2\pi}\left|e^{-it\psi_1\left(\sqrt{-\Delta}\right)} f_m   - e^{-it\psi_2\left(\sqrt{-\Delta}\right)} f_m \right|\right\|_{L^p(\X)} \lesssim 2^{-\frac{m\varepsilon}{2}} \sum_{j=1}^\infty \frac{1}{j!} \|g_{j,m}\|_{H^{\alpha_2}(\X)}\:.
\end{equation}
Recalling the definition of $g_{j,m}$ along with (\ref{transference_pf_eq3}), we note that
\begin{equation*}
\|g_{j,m}\|_{H^{\alpha_2}(\X)} \le C^j\:\|f_m\|_{H^{\alpha_2}(\X)}\:,
\end{equation*}
which upon plugging in (\ref{transference_pf_eq7}) gives,
\begin{eqnarray} \label{transference_pf_eq8}
\left\|\sup_{0\le t<2\pi}\left|e^{-it\psi_1\left(\sqrt{-\Delta}\right)} f_m   - e^{-it\psi_2\left(\sqrt{-\Delta}\right)} f_m \right|\right\|_{L^p(\X)} &\lesssim & 2^{-\frac{m\varepsilon}{2}} \|f_m\|_{H^{\alpha_2}(\X)}\left(\sum_{j=1}^\infty \frac{C^j}{j!}\right)\nonumber\\
& \lesssim &  2^{-\frac{m\varepsilon}{2}} \|f_m\|_{H^{\alpha_2}(\X)}\:.
\end{eqnarray}
Hence for each $m \ge m_0+1$, applying (\ref{transference_pf_eq1}) to $f_m$ and combining it with (\ref{transference_pf_eq8}), we get that
\begin{eqnarray*}
&&\left\|\sup_{0 \le t < 2\pi} \left|e^{-it\psi_2\left(\sqrt{-\Delta}\right)} f_m \right|\right\|_{L^p(\X)} \\
&\le & \left\|\sup_{0\le t<2\pi}\left|e^{-it\psi_1\left(\sqrt{-\Delta}\right)} f_m   - e^{-it\psi_2\left(\sqrt{-\Delta}\right)} f_m \right|\right\|_{L^p(\X)} + \left\|\sup_{0 \le t < 2\pi} \left|e^{-it\psi_1\left(\sqrt{-\Delta}\right)} f_m \right|\right\|_{L^p(\X)} \\
& \lesssim &  2^{-\frac{m\varepsilon}{2}} \|f_m\|_{H^{\alpha_2}(\X)} + \|f_m\|_{H^{\alpha_1}(\X)}\:.
\end{eqnarray*}
Then a further application of Lemma \ref{sobolev_comparison} gives
\begin{equation*}
\left\|\sup_{0 \le t < 2\pi} \left|e^{-it\psi_2\left(\sqrt{-\Delta}\right)} f_m \right|\right\|_{L^p(\X)} \lesssim 2^{-\frac{m\varepsilon}{2}} \|f_m\|_{H^{\alpha_2}(\X)} \le 2^{-\frac{m\varepsilon}{2}} \|f\|_{H^{\alpha_2}(\X)}\:,
\end{equation*}
which upon plugging in (\ref{transference_pf_eq4}) and combined with (\ref{transference_pf_eq5}) yields
\begin{eqnarray*}
&&\left\|\sup_{0 \le t < 2\pi} \left|e^{-it\psi_2\left(\sqrt{-\Delta}\right)} f \right|\right\|_{L^p(\X)}\\
 & \le & \left\|\sup_{0 \le t < 2\pi} \left|e^{-it\psi_2\left(\sqrt{-\Delta}\right)} f_{m_0} \right|\right\|_{L^p(\X)} + \sum_{m=m_0+1}^\infty \left\|\sup_{0 \le t < 2\pi} \left|e^{-it\psi_2\left(\sqrt{-\Delta}\right)} f_m \right|\right\|_{L^p(\X)} \\
&\lesssim & \left(1+\sum_{m=m_0+1}^\infty 2^{-\frac{m\varepsilon}{2}} \right)  \|f\|_{H^{\alpha_2}(\X)} \\
&\lesssim &  \|f\|_{H^{\alpha_2}(\X)}\:.
\end{eqnarray*}
As $\varepsilon>0$ and $K$-biinvariant $f \in C^\infty(\X)$ were arbitrarily chosen, (\ref{transference_pf_eq2}) follows. This completes the proof of Theorem \ref{transference_principle}.
\end{proof}

\end{document}
